\chapter{Our Glorious Past}
\section{Get Rolling}
\textbf{Here are a few seals from one of the earliest} civilisations \textbf{of the world.}
\vspace{0.1in}
\par
\textbf{What knowledge can we uncover about one of the earliest civilisations in the world by studying these ancient seals?}
\section{What to Know}
India's rich history is marked by significant contributions from the Indus Valley people and the Gupta rulers, who played pivotal roles in advancing architecture, trade, medicine, mathematics and science.
\section{\textbf{Learning Objectives}}
\begin{itemize}
    \item to describe the features of the Indus Valley towns
    \item to appereciate the skills of ancient builders and engineers of the Indus Valley and the Gupta period
    \item to discuss the contributions of the Indus Valley Civilisation and the Gupta period in the fields or trade, medicine, mathematics and science
    \item to analyse the importance of old inscriptions, religious texts, travel accounts, artefacts and clothes in reconstructing history
\end{itemize}
\section{I Explore}
\section{Indus Valley Civilisation}
The Indus Valley Civilisation was the first known civilisation of the Indian subcontinent spread over a large area, comprising parts of Pakistan and India. It is named after the river Indus, along the banks of which it flourished Harappa, Mohenjo-daro and Rakhigarhi were a few of the largest cities of this civilisation. The Indus Valley Civilisation existed roughly between 3000 BCE and 1500 BCE.
\vspace{0.5in}
\par
For thousands of years, the ruins of the Indus Valley Civilisation lay buried under the ground. In the 1920s, archaeologists began excavating near the Indus River and uncovered the vast city of Mohenjo-daro.
\vspace{0.1in}
\par
Over time, additional cities and numerous artefacts have been unearthed. It is only because of these artefacts discovered by archaeologists that we have been able to discover and understand the ancient Indus Valley Civilisation.
\vspace{2.2in}
\section{The Gupta Period}
The Gupta rulers ruled for over 200 years, from 320 CE to 550 CE, in the northern part of India.
\vspace{0.5in}
Pataliputra was the capital of the Gupta Empire. The Gupta rulers were great patrons of literature, science, art and architecture. Therefore, the Gupta period is also called he \textbf{Golden Age of India} because of tremendous growth in the fields of science, mathematics, medicine, art, literature and architecture.
\vspace{0.1in}
\par
The Indus Valley Civilisation and the Gupta Empire had a flourishingeconomic system. They made significant contributions to the fields of mathematics, astronomy and metallurgy. They were skilled builders and great engineers. Let us read about these people who lived in the first cities of the Indian subcontinent and those who lived during the Gupta period.
\section{Town Planning in the Indus Valley Civilisation}
The Indus Valley people lived in the first cities in India. They were excellent town planners. Their towns were planned and built better then many of the modern-day towns. All necessary services were arranged according to the convenience of the people. the cities were divided into two parts-upper town and lower town.
\subsection{Upper Town}
The Upper town was built on a raised platform and was higher than the lower town. It had a citadel (fortress). The public structures like the Great Bath, the Assembly Hall, etc. were found on the upper town of Moherjo-daro.
\subsection{\textbf{Public Buildings}}
Granaries were the most important public buildings in the citadel. They were close to the river so that grain could be transported easily. People depended upon these granaries for their food requirements, especially during famine days.
\vspace{0.1in}
\par
The citadel also had a great public bath-The Great Bath. This public bath was world famous for its excellent engineering. It was made watertight by using a layer of bitumen. It is believed that the general public used this bath on certain special festivals and ceremonies.
\vspace{0.1in}
\par
This glorious period of great buildings and engineering skills of the Indus Valley period remained unparalleled for a long, long time.
\subsection{Lower Town}
The lower town had residential buildings where the common people lived. The Harappans had advanced grid-based town planning and sanitation.
\subsection{\textbf{Houses}}
The houses in the lower town were built on both sides of the street. They were made of sun-dried or baked bricks. The entrance to the houses was from the side streets. this prevented crowding on the main streets. Each house had a well to draw water from and a kitchen and a bathroom with sloping floors.
\subsection{\textbf{Drainage}}
The sloping floors allowed water to move into the narrow drains built along the kitchens and the baths. These narrow drains carried the wastewater to the main underground drain built on the sides of the road. These drains carried the waste to large pits built outside the city. They were covered with stone slabs. They were inspected and cleaned at regular intervals by removing the stone slabs.
\vspace{0.4in}
\subsection{\textbf{Streets}}
The streets were constructed in straight lines and cut each other at right angles. The streets were built in this way so that they could be cleaned on their own by winds. They were also rounded off at corners to let bullock carts turn easily. They were made of burnt bricks and had a lighting system.
\section{Architecture under the Gupta Empire}
The Gupta kings were famous for building the first freestanding Hindu and Buddhist temples in India. They made square and flat-roofed rock temples. But later, they started using stones and bricks to make circular, rectangular and square temples with towers built over their roofs. This period marked the beginning of the magnificent stone structural temple architecture that used strong materials.
\vspace{0.1in}
\par
The Dashavarta Temple at Deogarh, who has towers, is a famous temple built by the Gupta kings. It is made of stone with sculpted wall panels on the ten \textit{avators} (incarnations) of Lord Vishnu.
\vspace{0.1in}
\par
Though the Gupta Empire lost its glory by the end of the sixth century, their contribution to Indian architecture will remain forever.
\vspace{1.2in}
\section{Scientist and Mathematicians}
Human have the potential to manipulate nature to their advantage. This desire has been responsible for the contributions to science, mathematics, metallurgy, medicine, surgery, trade and commerce. The people of the Indus Valley Civilisation and the Gupta period were great builders and engineers as well as scienists and mathematicians.
\subsection{Metal and Metallurgy}
The prople of the Indus Valley Civilisation knew about almost all metals except iron. They know how to make bronze by mixing copper and tin. They made many different objects like statues, tools and weapons of bronze. That is why it is also called a bronze-age civilisation.
\vspace{0.1in}
\par
In the Gupta period, people had a good knowledge of iron. The Iron Pillar in Delhi built by Chandragupta II is a to ancient India's knowledge and skill of metallurgy. This pillar has been out in the sun and rain for more then 1600 years but has not rusted till date.
\vspace{0.2in}
\par
I wonder know the Iron Pillar in Delhi, exposed the the Sun and rain for over 1600 years, has managed to resist rusting until today!
\vspace{0.2in}
\subsection{Archievements in Science and Mathematics}
The Indus people seem to have a good knowledge of mathematics. They used uniform weights made of special stones. these weights of different objects.
\vspace{0.1in}
\par
Aryabhatta was a great astronomer of the Gupta period. He gave astronomy a scientific viewpoint. He discovered zero. Without any advanced scientific instruments, he found out the exact time taken by Earth to rotate once around its own axis.
\vspace{0.1in}
\par
Warahamihira was another scientist of that time. He wrote a book of science called the \textit{Brihat-samhita.}
\vspace{0.4in}
\subsection{Medicine and Sugary}
Archaeologists have found that the Indus people used some plants, animal products and minerals as medicine. Denistry was also practised the Indus Valley Civilisation.
\vspace{0.1in}
\par
Knowledge of medicines greatly improved in the Gupta period. Takshila was a world-famous university renowned for science, especially medicine. Surgeons of that period could do eye, ear and nose surgeries. They could also treat and restore many injured or lost body parts. They knew how to set fractured and broken bones. There were hospitals to take care of the sick.
\vspace{0.1in}
The doctors practiced \textit{ayurveda} natural cures to treat patients. Sushruta was most likely the first surgeon to perform plastic surgery. In his book \textit{Sushrut-samhita}, he described surgical methods for bone setting, cataract removal and plastic sugary for repairing the nose, ears, lips, etc.
\section{Trade and Commerce}
Did you know that trade was well developed in the Indus Valley Civilisation and the Gupta Empire? The remains of a dockyard at Lothal in Gujurat prove that sea trade was flourishing in India as early as the Indus Valley Civilisation.
\vspace{0.1in}
\par
The presence of Harappan seals and artefacts in the Mesopotamian and Egyptian civilisations in Harappan sites is evidence of trade between these ancient civilisations.
\vspace{0.1in}
\par
Traders of different towns and cities of the Indus valley exchanged cotton, wheat, some other food crops, pots and animal products, tools weapons and ornaments with each other. They bought copper, tin, gold, silver and other costly metals from Mesopotamian (modern Iraq), Persia (modern Iran), Oman and Afghanistan.
\vspace{0.5in}
\par
On the other hand, traders of the Gupta period sold cloth, food grains, spices, salt and precious stones to different parts of India and the world. They bought gold, silver, tin, lead, silk and horses from the countries in West Asia, Europe, Southeast Asia like Java and Sumatra (both islands in modern Indonesia), Cambodia and China. They sold pearls, precious stones, silk, perfumes, spices, indigo, drugs, coconuts and herbs to these countries.
\vspace{0.1in}
\par
A strong economy in the ancient period must have led to an organised social structure and a good standard of living. India was technologically, culturally, socially advanced and economically self-sufficient. Hence, our glorious past is a matter of pride for all of us. It inspires us to work towards taking India to greater heights!
\vspace{2in}
\section{Now I Know}
\begin{itemize}
    \item The Indus Valley Civilisation was the first known civilisation of the INduan subcontinent. It was located near the river Indus.
    \item The Indus Valley Civilisation existed roughly between 3000 BCE and 1500 BCE.
    \item The Gupta rulers ruled for over 200 years, from 320 CE to 550 CE, in the northern part of India.
    \item The Gupta period is also called the golden age of India bevause of tremendous growth in the fields of science, mathematics, medicine, art, literature and architecture during this period.
    \item The Indus valley people were excellent town planners.
    \item The houses in the Indus Valley Civilisation were made of sun-dried or baked bricks.
    \item The Gupta kings were famous for building the first freestanding Hindu and Buddhist temples in India.
    \item The Iron Pillar in Delhi built by Chandragupta II is a testament to ancient India's knowledge and skill of metallurgy.
    \item Aryabhatta was a great astronomer of the Gupta period. He discovered zero. Without any advanced scientific instruments, he found out the exact time taken by Earth to rotate once around its own axis.
    \item Both internal and foreign trade were well developed in the Indus Valley Civilisation and the Gupta Empire.
\end{itemize}